% Two update paths, and what the second one is worth.
\begin{tikzpicture}[font=\sffamily\scriptsize, x=1cm, y=1cm]

\node[chlabel, anchor=west] at (-6.8,3.6) {two ways to change this node's firmware, and only one of them had to be written};

% ================================================================ the ROM path
\node[layerlabel] at (-6.8,3.1) {already in the silicon, and it speaks this bus};
\node[ext, text width=24mm] (m1) at (-5.2,2.2) {the master};
\node[hw, text width=26mm] (rom) at (-1.6,2.2) {the ROM bootloader\\two pins, two bit rates};
\node[hw, text width=24mm] (fl1) at (1.8,2.2) {flash};
\draw[flow] (m1) -- (rom);
\draw[flow] (rom) -- (fl1);
\node[chlabel, anchor=west, text width=36mm] at (3.4,2.2)
  {and its detection chain polls the bus \textbf{first}, before the serial and
   USB routes};

% ================================================================ what it does not do
\node[regcell, fill=red!18, minimum width=114mm, text width=112mm, inner sep=3pt,
      anchor=north west, align=left] at (-6.8,1.4)
  {\textbf{What it does not do:} report a version, refuse a corrupt image, carry
   a node identifier of its own, answer a fleet sweep, or let the bit rate be
   anything other than the two it fixes. \textbf{Those five absences are the
   entire justification for the rest of this chapter}, and a custom loader that
   does not beat something already in the silicon is not worth its flash.};

% ================================================================ the custom path
\node[layerlabel] at (-6.8,-0.2) {and the one this chapter writes};
\node[ext, text width=24mm] (m2) at (-5.2,-1.1) {the master};
\node[fw, text width=26mm] (bl) at (-1.6,-1.1) {the loader\\a polled superloop};
\node[hw, text width=24mm] (fl2) at (1.8,-1.1) {flash};
\draw[flow] (m2) -- (bl);
\node[chlabel, anchor=north] at (-3.45,-1.6) {segmented transfer};
\draw[flow] (bl) -- (fl2);
\node[laterblk, text width=24mm] (app) at (1.8,-2.4) {the application\\a task set};
\draw[flow] (bl) -- node[ctrllabel, right]{checksum, then jump} (app);
\draw[recoveredge] (app) to[bend left=35] node[buslabel]{a RAM keyword, then reset} (bl);

% ================================================================ the requirements
\node[regcell, fill=hwcol, minimum width=54mm, text width=52mm, inner sep=3pt,
      anchor=north west, align=left] at (-6.8,-3.4)
  {\textbf{Requirements, from two primary documents read in full:} it always
   executes first; it always computes the application checksum and never runs a
   mismatch; it accepts exactly two identifiers; every received block is bounds
   checked; and its own sectors are protected.};

\node[regcell, fill=usrcol, minimum width=54mm, text width=52mm, inner sep=3pt,
      anchor=north west, align=left] at (0.2,-3.4)
  {\textbf{And one simplification that comes from this part:} the vector table
   offset register is programmable here, so the loader points it at the
   application before jumping. Designs on parts with a fixed vector location
   have to mirror the table and route every vector through a decider.};

\node[umlnote, text width=114mm, anchor=north west] at (-6.8,-6.4)
  {The bounded wait before the application starts is a cost paid at every single
   power-up, in exchange for the node always being reachable afterwards. That is
   a trade rather than an oversight, it is stated in the budget table with a
   number beside it, and a design that does not state it has usually not noticed
   it.};
\end{tikzpicture}
